\subsection{Setup and Formal Definition}
Let $x[n]$ be a discrete-time signal with $n \in \mathbb{Z}$, sampled at frequency $f_s$.
Patch-based tokenisation extracts a sequence of tokens by sliding a window of size $P$
samples with stride $S$:
\begin{equation}
    \mathbf{p}_k
    = \bigl(x[kS],\; x[kS+1],\; \ldots,\; x[kS+P-1]\bigr)
    \in \mathbb{R}^P,
    \qquad k = 0, 1, 2, \ldots
\end{equation}

\begin{definition}[Patch Tokenisation Operator]
    The map $\mathcal{T}: \ell^\infty(\mathbb{Z}) \to (\mathbb{R}^P)^{\mathbb{N}}$ defined
    by~(1) is called the \emph{patch tokenisation operator}.  It is linear: collecting $K$
    patches yields
    \begin{equation}
        \mathbf{P} = \mathcal{T}\{x\} = \mathbf{W} \cdot \mathbf{x},
    \end{equation}
    where $\mathbf{W} \in \{0,1\}^{KP \times N}$ is a binary selection matrix whose $k$-th
    block of $P$ rows picks the $P$ consecutive samples starting at index $kS$.
\end{definition}

\subsection{Condition for Two Consecutive Patches to Be Identical}

\begin{proposition}
    Patches $\mathbf{p}_k$ and $\mathbf{p}_{k+1}$ are identical if and only if
    \begin{equation}
        x[n] = x[n + S] \quad \forall\, n \in \{kS,\, kS+1,\, \ldots,\, kS+P-1\}.
    \end{equation}
    Globally, if the signal is periodic with period $T_0$ (in samples), identical
    consecutive patches occur whenever
    \begin{equation}
        \boxed{S = m \cdot T_0, \qquad m \in \mathbb{Z}^+.}
    \end{equation}
    Equivalently, in terms of signal frequency $f$ and sampling frequency $f_s$:
    \begin{equation}
        f = m \cdot \frac{f_s}{S}, \qquad m \in \mathbb{Z}^+.
    \end{equation}
\end{proposition}

\begin{proof}
    $\mathbf{p}_k = \mathbf{p}_{k+1}$ component-wise requires
    $x[kS + j] = x[(k+1)S + j]$ for all $j \in \{0,\ldots,P-1\}$,
    i.e.\ $x[n] = x[n+S]$ on that window.
    For a periodic signal with $T_0 = f_s/f$ samples, this holds whenever $S$ is an
    integer multiple of $T_0$, giving~(4).
\end{proof}

\subsection{Loss of Observability}

When $\mathbf{p}_k = \mathbf{p}_{k+1}$ for all $k$, the token sequence is constant — the
model sees the same vector repeated.  The \emph{rank} of the token matrix $\mathbf{P}$
collapses to 1.  Consequently:

\begin{itemize}
    \item Temporal variation is invisible to the model.
    \item Any two signals differing only by a phase shift of $S$ samples produce the same
          token sequence: they are \textbf{observationally equivalent}.
    \item Even when patches are not fully identical but merely highly correlated
          ($S \approx mT_0$), the effective information in the token sequence is reduced
          (\emph{soft} loss of observability).
\end{itemize}